\documentclass{article}

\usepackage{geometry}
\usepackage{lipsum}
\usepackage{tcolorbox}
\usepackage{mdframed}
\usepackage{amsmath}
\usepackage{texmaths-colors}
\usepackage{dashrule}

\newgeometry{
        lmargin=1.5cm,
        rmargin=1.5cm,
        top=1.5cm,
        footskip=0pt
    }%

\mdfdefinestyle{byheart}{
        backgroundcolor=white!80!black,
        linecolor=white!20!black,
        fontcolor=red!0!black,
        roundcorner=5pt,
        leftmargin=\parindent,
        innertopmargin=8pt,
        innerbottommargin=8pt,
    }
\mdfdefinestyle{simple}{
        backgroundcolor=white!90!black,
        linecolor=white!10!black,
        fontcolor=red!0!black,
        roundcorner=5pt,
        leftmargin=\parindent
        }

\def\titre#1{\begin{mdframed}[style=byheart]
    \centering\huge #1
\end{mdframed}}

\newlength{\boxwidth}

\long\def\methode#1#2{\settowidth{\boxwidth}{\Large #1}
\advance\boxwidth20pt\vskip20pt\noindent
\begin{minipage}[b]{\boxwidth}
\begin{mdframed}[style=simple]
    {\Large #1}
\end{mdframed}
\end{minipage}\hskip10pt
{\Large\color{vscode-crimson} #2}\vskip20pt}

\def\calcstretch{2.5}
\def\mulby#1#2{#1{\color{vscode-blue}\times#2}}
\def\newcalc{\hdashrule{.7\linewidth}{.5pt}{2pt}\vskip20pt}

\begin{document}
\pagestyle{empty}\large
\titre{Fiche méthode : additions et soustraction de fractions}
\methode{Les dénominateurs sont identiques}{On ajoute les numérateurs.}
\begingroup\renewcommand{\arraystretch}{\calcstretch}
    $\begin{array}{cccc}
        &\dfrac23&+&\dfrac53\\
        =&&\dfrac73&
    \end{array}$
\endgroup
\methode{Les dénominateurs sont différents}{On cherche un multiple commun.}
\begin{minipage}{.4\linewidth}\renewcommand{\arraystretch}{\calcstretch}
    $\begin{array}{cccc}
        &\dfrac2{21}&+&\dfrac57\\
        =&\dfrac2{21}&+&\dfrac{\mulby53}{\mulby73}\\
        =&\dfrac2{21}&+&\dfrac{15}{21}
    \end{array}$
\end{minipage}
\begin{minipage}{.6\linewidth}
    \color{vscode-blue}On remarque que $7\times3=21$.
\end{minipage}\vskip20pt\newcalc
\begin{minipage}{.4\linewidth}\renewcommand{\arraystretch}{\calcstretch}
    $\begin{array}{cccc}
        &\dfrac5{6}&+&\dfrac{10}9\\
        =&\dfrac{\mulby{5}3}{\mulby63}&+&\dfrac{\mulby{10}2}{\mulby{9}2}
    \end{array}$
\end{minipage}
\begin{minipage}{.6\linewidth}
    \color{vscode-blue}On remarque que 18 est un multiple de 6 et 9.
\end{minipage}\vskip20pt\newcalc
\begin{minipage}{.4\linewidth}\renewcommand{\arraystretch}{\calcstretch}
    $\begin{array}{cccc}
        &\dfrac{13}{11}&+&\dfrac37\\
        =&\dfrac{\mulby{13}7}{\mulby{11}7}&+&\dfrac{\mulby3{11}}{\mulby7{11}}
    \end{array}$
\end{minipage}
\begin{minipage}{.6\linewidth}
    \color{vscode-blue}Si on ne trouve pas un multiple commun,
    il suffit de multiplier les dénominateurs~:$$11\times7=7\times11=77$$
\end{minipage}
\end{document}